\begin{center}
  \Large
  \textbf{KATA PENGANTAR}
\end{center}

\addcontentsline{toc}{chapter}{KATA PENGANTAR}

\vspace{2ex}

% Ubah paragraf-paragraf berikut dengan isi dari kata pengantar

Puji dan syukur kehadirat Tuhan Yang Maha Esa, atas segala rahmat dan karunia-Nya, sehingga penulis dapat menyelesaikan penelitian ini yang berjudul \tatitle

Penelitian ini disusun dalam rangka pemenuhan Tugas Akhir sebagai syarat kelulusan Mahasiswa ITS. Oleh karena itu, penulis mengucapkan banyak terima kasih kepada:

\begin{enumerate}[nolistsep]
  \item Bapak Fuad dan Bapak Ketut selaku pembimbing penelitian ini yang telah memberikan arahan, masukan, dan  motivasi selama penyusunan tugas akhir ini.
  \item Orang tua dan keluarga tercinta yang selalu memberikan doa serta dukungan secara moral maupun material.
  \item Bapak dan ibu dosen pengajar Departemen Teknik Komputer yang telah memberikan ilmu dan pengalaman selama masa kuliah.
  \item Bapak dan ibu tenaga kependidikan Departemen Teknik Komputer yang telah membantu dalam administrasi dan sumber informasi selama masa penulisan ini.
  \item Aurellia Edelweis yang telah memberikan dukungan mental dan telah menemani selama penyusunan tugas akhir ini.
  \item M. Raidan Arsyal Y.K. yang telah menemani hingga akhir penyusunan tugas akhir.
  \item Teman-teman Teknik Komputer angkatan 2021, khususnya Fardian A. Zidane, Ezekiel Walfred E.P.N., M. Rafi Irsyad S., dan Ardy Pratama S. yang telah menemani dan memberi motivasi selama masa kuliah.
  \item Teman-teman Teknik Komputer yang telah memberikan bantuan dan motivasi untuk menyelesaikan tugas akhir ini. 
  \item Para kontributor yang telah berkontribusi dalam upaya pelestarian satwa, terutama penyu laut di dunia.
  \item Para pengidap gangguan ADHD yang telah memberikan bantuan saran dan motivasi dalam menjalani hidup.
\end{enumerate}

Akhir kata, semoga penelitian ini dapat memberikan manfaat kepada banyak pihak. Penulis menyadari jika tugas akhir ini masih jauh dari kata sempurna. Untuk itu penulis mengharapkan saran dan kritik yang bersifat membangun untuk dapat menuai hasi yang lebih baik lagi.

\begin{flushright}
  \begin{tabular}[b]{c}
    \place{}, \MONTH{} \the\year{} \\
    \\
    \\
    \\
    \\
    \name{}
  \end{tabular}
\end{flushright}
